%%=============================================================================
%% Samenvatting
%%=============================================================================

% De "abstract" of samenvatting is een kernachtige (~ 1 blz. voor een
% thesis) synthese van het document.
%
% Een goede abstract biedt een kernachtig antwoord op volgende vragen:
%
% 1. Waarover gaat de bachelorproef?
% 2. Waarom heb je er over geschreven?
% 3. Hoe heb je het onderzoek uitgevoerd?
% 4. Wat waren de resultaten? Wat blijkt uit je onderzoek?
% 5. Wat betekenen je resultaten? Wat is de relevantie voor het werkveld?
%
% Daarom bestaat een abstract uit volgende componenten:
%
% - inleiding + kaderen thema
% - probleemstelling
% - (centrale) onderzoeksvraag
% - onderzoeksdoelstelling
% - methodologie
% - resultaten (beperk tot de belangrijkste, relevant voor de onderzoeksvraag)
% - conclusies, aanbevelingen, beperkingen
%
% LET OP! Een samenvatting is GEEN voorwoord!

%%---------- Nederlandse samenvatting -----------------------------------------

\IfLanguageName{english}{%
\selectlanguage{dutch}
\chapter*{Samenvatting}
\lipsum[1-4]
\selectlanguage{english}
}{}

%%---------- Samenvatting -----------------------------------------------------
% De samenvatting in de hoofdtaal van het document

\chapter*{\IfLanguageName{dutch}{Samenvatting}{Abstract}}


De toenemende adoptie van SaaS-platformen en DevOps-praktijken brengt heel wat uitdagingen mee voor security monitoring. Organisaties werken met gedecentraliseerde infrastructuren waarin logs, alerts en events verspreid zitten over meerdere tools zonder onderlinge koppeling. Voor Devinity, een SaaS-gericht bedrijf, vertaalt zich dat in een situatie waarbij security-events niet systematisch worden gedetecteerd, incident response voornamelijk reactief verloopt en meetpunten om de effectiviteit van de aanpak te beoordelen volledig ontbreken.

De centrale onderzoeksvraag van deze bachelorproef luidt: hoe kan een geautomatiseerd systeem voor security monitoring en incident response worden ontworpen om securitydata binnen de SaaS-omgeving van Devinity effectief te centraliseren, analyseren en rapporteren in een DevOps-infrastructuur? Het doel is een werkend proof-of-concept te bouwen dat de haalbaarheid van zo'n aanpak aantoont binnen de bestaande toolstack van het bedrijf, zonder de infrastructuur te vervangen.

Het onderzoek verloopt in vier fasen. In de probleemanalyse worden de concrete knelpunten in kaart gebracht via observaties en gesprekken met het team. Op basis daarvan wordt een gelaagde architectuur ontworpen. Make werkt als ingestie- en filterlaag, Azure Log Analytics als centraal opslag- en analyseplatform, Azure Monitor Workbooks als visualisatielaag en monday.com in combinatie met Outlook voor verbeterde incident response. Vervolgens wordt deze architectuur uitgewerkt tot een werkend proof-of-concept dat de Azure WAF als primaire databron gebruikt, via Azure Monitor Alert Rules, KQL-queries, een Data Collection Rule en een Make-scenario met routeringslogica op basis van de ernst van een event. In de vierde fase worden de resultaten geëvalueerd aan de hand van drie KPI's.

Na de evaluatie ligt de gemiddelde detectietijd op 187 seconden, consistent binnen het 5-minutenvenster van de Alert Rules. De gemiddelde totale responstijd van trigger tot notificatie of workitem is nog maar 23,9 seconden, waarbij Make-verwerking gemiddeld slechts 2,8 seconden inneemt. Over vijf testrondes werd een filterratio van 98,5\% gemeten. Van de in totaal 3\,857 ruwe events bereikten er slechts 59 de centrale opslag, resulteerden er 8 in een workitem en 2 in een e-mailnotificatie. Die hoge ratio is te verklaren door het beperkte verkeersvolume in de QA-omgeving, maar de filterlogica functioneert zoals ontworpen.

De resultaten tonen aan dat een betere aanpak van security monitoring haalbaar is zonder nieuwe tools te gebruiken of de bestaande infrastructuur te vervangen. De meerwaarde voor het werkveld is dat de architectuur generiek is opgebouwd en daardoor ook bruikbaar is in andere SaaS- en DevOps-omgevingen met gelijkaardige uitdagingen. De voornaamste beperking is dat de validatie beperkt bleef tot de QA-omgeving met gesimuleerde testdata. Cross-source correlatie en productievalidatie blijven aanbevelingen voor een volgende fase.
